% !PRINTED 457
\weberpage{457}
If we denote the coefficients of $\varphi(x)$, which, as we have seen, are quantities in $\Omega$, by $c_0, c_1 \ldots c_{n-1}$, then every quantity $\Theta$ in $\Omega(\alpha)$ can accordingly be represented in the form:

\[
\Theta = c_0 + c_1\alpha + c_2\alpha^2 + \cdots + c_{n-1}\alpha^{n-1}.
\]

This representation is possible in only \gesperrt{one way}; for if at the same time

\weq{6}{\Theta = c_0' + c_1'\alpha + c_2'\alpha^2 + \cdots + c_{n-1}'\alpha^{n-1},}

then the function of degree $(n-1)$

\[
c_0 - c_0' + (c_1 - c_1')x + (c_2 - c_2')x^2 + \cdots + (c_{n-1} - c_{n-1}')x^{n-1}
\]

vanishes for $x = \alpha$. But by Theorem III of the preceding section this is possible only if

\[
c_0 = c_0', \quad c_1 = c_1', \ldots c_{n-1} = c_{n-1}'
\]

holds.

If we regard the quantities of the field $\Omega$ as the given ones, its quantities as rational, and the field $\Omega(\alpha)$ as the one derived from them, then we also call a rational function in $\Omega$ briefly a \gesperrt{rational function}, and can accordingly also say:

\gesperrt{The field $\Omega(\alpha)$ consists of all rational functions of $\alpha$.}

\begin{center}
\S.\,143.

\gesperrt{Simultaneous adjunction of several algebraic quantities.}
\end{center}

To a more exact investigation of the field $\Omega(\alpha)$ we prefix a general theorem, very simple but yet rich in consequences.

1. \gesperrt{If $\Phi_1(x,y,z\ldots)$, $\Phi_2(x,y,z\ldots)$, $\Phi_3(x,y,z\ldots)\ldots$ are integral rational functions of the variables $x,y,z\ldots$ with arbitrary coefficients, which do not all vanish at once in any of the functions, then one can substitute for the variables, in infinitely many ways, such rational numerical values that none of the functions $\Phi_1, \Phi_2, \Phi_3\ldots$ vanishes.}
